\documentclass{paper}

\usepackage{fancyhdr}
\usepackage[a4paper, total={170mm,257mm}, left=20mm, top=20mm]{geometry}

\begin{document}

\section*{Mise en cohérence des objectifs du TIPE (MCOT) - GIL Dorian}

\section{Intro}
\begin{itemize}
    \item \textbf{Titre: }Application de la méthode des tableaux en logique propositionelle et en logique linéaire temporelle
    \item \textbf{Ancrage: }La méthode des tableaux est un algorithme qui renvoie la satisfiabilité d'une formule. Nous étudierons une variante où le tableau aura la forme d'un arbre [7]. Une formule est insatisfaisable ssi toutes les branches ont des cycles (c'est à dire $\phi$ et $\lnot\phi$ sur une même branche)
    \item \textbf{Motivations: }Etant interessé par mes cours de logiques, et après des recherches, j'ai découvert la méthode des tableaux. Après plus de recherches notamment dans ces applications, j'ai découvert et je me suis intéressé à la vérification de modèle, d'où l'étude de la méthode des tableaux dans le cadre de la vérification de modèle.
\end{itemize}

\section{Positionnements thématiques et mots-clés}
\subsection{Thème}
\begin{enumerate}
\item Informatique - Théorique
\item Informatique - Pratique
\item Mathématiques - Autres (Logiques classiques et non-classiques ; Algorithmique)
\end{enumerate}

\subsection{Mots clés - Key words}
\begin{enumerate}
\item Méthode des tableaux // Method of Analytic Tableaux
\item Logique propositionnelle // Propositional Logic
\item Logique linéaire temporelle // Linear Temporal Logic
\item Satisfiabilité // Satisfiability
\item Vérification de modèle // Model checking
\end{enumerate}

\section{Bibliographie commentée (493/650 mots)}

Le theorème de Cook nous indique que le problème de décision de la satisfiabilité d'une formule est NP-complet, c'est ainsi que la question de la 
satisfiabilité en temps polynomial est devenue question à un million de dollars, la résolution de ce colossal problème impliquerait en effet P = NP.  
Donc une grande problématique dans le domaine de la logique est la recherche d'un algorithme permettant de résoudre le problème de satisfiabilité en 
temps polynomial.

Ainsi, dans le cadre de la logique propositionnelle, la méthode des tableaux consiste à construire un  arbre avec la formule à la racine, et à utiliser 
des règles pour développer ou créer des branches. On regarde ensuite s'il y a des contradictions/cycles dans toutes les branches. Si c'est le cas, la 
formule est insatisfaisable. La méthode des tableaux résout le problème SAT en temps exponentiel en la taille de la formule étudiée. Au délà de la 
résolution de P=NP, trouver des optimisations permet in fine de résoudre en temps raisonnable d'autres problèmes pouvant se réduire à une recherche 
de satisfaibilité ou d'accélerer la preuve automatique de certaines familles de formules logiques.

Néanmoins, ces problèmes se transfèrent sur d'autres types de logiques, en particulier la logique linéaire temporelle, qui est une logique dans laquelle
on rajoute une notion de temporalité representé de manière discrète par une suite infinie d'évaluations de vérité des variables. Ainsi les variables 
peuvent être valuer differement dans le temps. La méthode des tableaux est un algorithme qui est particulièrement utilisé dans la logique linéaire 
temporelle dans le cadre de la vérification de modèle.

La vérification de modèle est un problème dans lequel on souhaite vérifier si le modèle d'un système satisfait une propriété. Il existe des algorithmes 
résolvant ce problème dans le cadre la logique temporelle. Un de ces algorithmes de vérification automatique en logique linéaire temporelle utilise
la méthode des tableaux.

C'est ainsi que, dans un monde toujours plus dépendant des différentes technologies dont la création repose sur des algorithmes, on souhaite vérifier le
bon fonctionnement de nos programmes informatiques par rapport aux cahiers des charges associés à un produit. La vérification de modèle et les différents
algorithmes qu'elle offre permet de répondre à ce problème. D'où l'enjeux majeur de la méthode des tableaux dans le cadre de la logique linéaire temporelle.

Naturellement, comme tout modèle voulant representé la realité, nous pouvons nous interogger sur la manière dont on doit procéder, pour ainsi après 
utilisation d'un algorithme de vérification de modèle prouvé rigoureusement le bon fonctionnement de notre système.

\section{Problématique retenue}
Comment utiliser la méthode des tableaux pour la résolution du problème de décision de la satisfiabilité d'une formule logique et de la vérification de modèle ?

\section{Objectifs du TIPE}
Les objectifs du TIPE sont d'étudier les différentes utilisations de la méthode des tableaux dans différentes logiques. Nous nous intéresserons en particulier :

Premièrement, dans la logique propositionnelle avec une implémentation personelle de la méthode des tableaux avec en particulier une étude de la satisfiabilité d'une 
famille de formules, dont on implémentera l'algorithme suivant la méthode des tableaux qui résout la satisfiabilité en temps optimisé.

Deuxièmement, dans la logique linéaire temporelle avec l'étude d'une formule logique modélisant le fonctionnement attendu d'un objet technologique, on
soulignera en particulier les hypothèses faites pendant la modélisation et l'intérêt de la méthode des tableaux après cette phase de modélisation, en utilisant
une implémentation de celle ci.

\section{Liste de références bibliographiques} 
\begin{enumerate} 
    \item Pierre Le Barbenchon, Sophie Pichinat, François Schwarzentruber ; Logique : fondements et applications ; Dunod ; ISBN : 978-2-10-082158-7 
    \item Chiara Ghidini (Università di Trento) ; Mathematical Logic: Tableaux Reasoning for Propositional Logic ; https://dit.unitn.it/\verb|~|ldkr/ml2015/slides/PLtableau.pdf 
    \item Garey M.R., Johnson D.S. ; Computers and Intractability: A Guide to the Theory of NP-Completeness, 1 ed.; San Francisco: W. H. Freeman and Company, 1979 
    \item The tableau method for temporal logic: an overview ; Pierre Wolper ;01 Jan 1985 - Logique Et Analyse - Vol. 28, pp 119-136 
    \item Principles of Model Checking ; Christel Baier et Joost-Pieter Katoen ; Chap 5 - 9780262026499 - April 25, 2008 - The MIT Press 
    \item Simple On-the-fly Automatic Verification of Linear Temporal Logic ; Rob Gerth, Doron Peled, Moshe Y. Vardi, Pierre Wolper ; https://orbi.uliege.be/bitstream/2268/116644/1/GPVW95-pstv.pdf 
    \item A traditional tree-style tableau for LTL ; Mark Reynolds ; https://arxiv.org/pdf/1604.03962 - The University of Western Australia ;
\end{enumerate}

\end{document}